% Section 4.2 -- Decoupling. Plain, direct, number-backed prose (RAG Results register).
\subsection{Decoupling reasoning from the gradient}
\label{sec:decoupling}

The rate target alone is not enough, because a chain of thought before the answer lets the same reward
make calibration worse. Applied to the full completion, the reward of \cref{eq:reward} drives the
held-out Brier score from the base model's $0.25$ to $0.34$ and the calibration error from $0.19$ to
$0.30$ (\cref{fig:trajectory}). The policy gradient passes through the reasoning tokens, so raising the
reward on the final probability rewrites the reasoning before it into pseudo-quantitative arguments for
ever more extreme numbers. This is the accuracy-calibration gradient conflict reported for
verifiable-reward training \citep{dcpo}.

Two changes keep the reward and stop the gradient from reaching the reasoning. The first, the direct
model, drops the chain of thought and emits the probability directly. With no reasoning tokens to
corrupt, the reward calibrates the model, and the held-out Brier score reaches the teacher's level
within fifty steps and stays there. The second, the masked model, keeps the chain of thought but masks
the gradient. We narrow the completion mask, which selects the tokens that receive advantage, to the
final ``Probability: NN\%'' span, so the reasoning is still sampled and still conditions the answer but
gets no gradient. Masking the gradient
needs the KL penalty turned off, because on the few tokens of the answer span the gradient is
concentrated enough to make the penalty's estimator overflow and the loss diverge, and the masked
reasoning no longer needs the reference-model anchor. With the mask in place, the masked model
calibrates as well as the direct model.

\begin{figure*}[t]
\centering
\includegraphics[width=\linewidth]{F2_trajectory.pdf}
\caption{Held-out calibration during training (fixed $n=128$). Applied to the full completion, the
reward of \cref{eq:reward} drives a chain-of-thought policy away from calibration on both Brier (a) and
ECE (b). Removing the reasoning from the gradient, by direct prediction or by masking the gradient to
the answer span, recovers calibration to the teacher's level. The full-completion and masked runs share
the same reward, prompt, and base model, and differ only in the gradient mask.}
\label{fig:trajectory}
\end{figure*}

The masked model is slightly less accurate than the direct model. On held-out 2024 games the direct
model reaches a Brier of $0.144$ and the masked model $0.152$, a gap of $0.008$ by a paired bootstrap
over plays (95\% CI $[0.006, 0.010]$). The gap is small, and far below the base-to-teacher gap that the
reward closes. Both trained models trail the market on Brier, but the shortfall is in resolution, not
calibration: the direct model's calibration error of $0.029$ matches the market's $0.027$, and its Brier
is higher only because the static prompt cannot see the live information the market uses to separate
states more sharply (\cref{sec:results}).

\begin{table}[t]
\centering
\setlength{\tabcolsep}{4.5pt}
% paired bootstrap, per-play Brier difference (negative => row-A better). $\ast$=95% CI excludes 0.
\begin{tabular}{lcc}
\toprule
Paired comparison (Brier) & $\Delta$ & 95\% CI \\
\midrule
Direct $-$ Masked-CoT & -0.0079$^{\ast}$ & {\scriptsize[-0.0101,\,-0.0057]} \\
Direct $-$ Market & +0.0088$^{\ast}$ & {\scriptsize[0.0068,\,0.0107]} \\
Masked-CoT $-$ Market & +0.0167$^{\ast}$ & {\scriptsize[0.0140,\,0.0194]} \\
Direct $-$ DeepSeek-V4 & +0.0005 & {\scriptsize[-0.0019,\,0.0028]} \\
DeepSeek-V4 $-$ Market & +0.0083$^{\ast}$ & {\scriptsize[0.0064,\,0.0102]} \\
\bottomrule
\end{tabular}

\caption{Paired bootstrap over held-out 2024 plays ($n=5185$, $10^4$ resamples). Masking the gradient
to the answer span costs a small but significant amount of Brier relative to discarding the reasoning;
both decoupled models trail the market by a similar, significant margin that \cref{sec:results}
attributes to resolution rather than calibration.}
\label{tab:decoupling}
\end{table}

The masked model keeps its reasoning coherent. A blinded judge read each completion and decided whether
the stated probability follows from the reasoning, and the rate of inconsistent completions falls from
$22.4\%$ for the base model to $4.4\%$ after masked training. Full-completion reinforcement learning
leaves the base rate unchanged, so the gain comes from the masked objective and not from reinforcement
learning in general, and a control that applies the masked prompt without training accounts for only
part of it. The base model writes arithmetic that does not sum and then reports an unrelated number, or
rates a clear favorite below an even chance, while the masked model reads the live state and states a
probability that follows from it (\cref{fig:reasoning-examples}). A probability lookup cannot do this,
so the masked model is the better choice when a forecast must be explained, at the cost of slightly
lower sharpness than direct prediction.

\begin{figure*}[t]
\centering
\includegraphics[width=\linewidth]{F5_faithfulness.pdf}
\caption{Rate at which the stated probability does not follow from the model's own chain of thought, on
a blinded-judge sample of held-out plays (95\% bootstrap intervals). Full-completion chain-of-thought
reinforcement learning leaves the base rate unchanged; masked training reduces it sharply, beyond what
the masked prompt format alone accounts for.}
\label{fig:faithfulness}
\end{figure*}

% E1 -- faithful vs. corrupted reasoning. Curated verbatim excerpts from the blinded-judge
% dump (paper/examples/E1_candidates.txt; eval/results/reasoning_judged_eval.jsonl, eval 2023).
% Excerpts trimmed to the decisive lines; [...] marks elision; probabilities are the models' own.
% Requires \usepackage{xcolor}; uses \ourdirect/\ourmasked colors if defined, else falls back.
\providecommand{\statehead}[1]{\smallskip\noindent\textbf{#1}\par\smallskip}

\begin{figure*}[t]
\centering
\small
\begin{minipage}{0.99\textwidth}

\statehead{Example 1.\quad Q4, 2:21 left. WAS trailing by 12, 2nd \& 6 at the NYG 8-yard line.
Pregame line: WAS favored by 7.5. Empirical rate $\hat{p}=0.07$; WAS lost.}

\begin{tabular}{@{}p{0.085\textwidth}@{\hskip 6pt}p{0.80\textwidth}@{\hskip 6pt}r@{}}
\textit{Base (CoT)} &
\emph{``[\,\dots] Touchdown and Extra Point: $0.8\times0.95=0.76$ or 76\%. Touchdown and
Two-Point Conversion: $0.8\times0.4=0.32$ or 32\%. Total Probability: $0.76+0.32=1.08$, but
since we can't exceed 100\%, we cap it at 100\%. [\,\dots] approximately 32\%.''} & $0.32$ \\[3pt]
\textit{Masked-CoT (ours)} &
\emph{``[\,\dots] a 12-point deficit with only 2 minutes and 21 seconds left [\,\dots] makes a
comeback extremely challenging. Probability: 8\%.''} & $0.08$ \\
\end{tabular}

\medskip\hrule\medskip

\statehead{Example 2.\quad Q4, 7:57 left. CHI leading by 3, 3rd \& 6 at the CAR 42-yard line.
Pregame line: CHI favored by 3. Empirical rate $\hat{p}=0.80$; CHI won.}

\begin{tabular}{@{}p{0.085\textwidth}@{\hskip 6pt}p{0.80\textwidth}@{\hskip 6pt}r@{}}
\textit{Base (CoT)} &
\emph{``[\,\dots] Base Probability (Favored by 3): 60--70\%. Adjustment for Poor Performance:
$-10$ to $-20\%$. Final Probability: 40--50\%.''} \ (a leading, favored team, with the stated
value at the bottom of its own adjusted range) & $0.40$ \\[3pt]
\textit{Masked-CoT (ours)} &
\emph{``[\,\dots] CHI is in a strong position to maintain their lead, as they are favored and
have the ball [\,\dots] ample time to either score a touchdown or at least a field goal.
Probability: 85\%.''} & $0.85$ \\
\end{tabular}

\end{minipage}
\caption{Reasoning before and after gradient-masked CoT training, on identical game states
(blinded-judge sample, 2023). In the base completions the stated probability does not follow from
the text: an addition that sums to 108\% reports one of its branches, and a leading favorite is
rated below an even chance. The masked-CoT model reads the live state and states a probability
consistent with it, close to the empirical rate $\hat{p}$. Inconsistency of this kind falls from
22.4\% of base completions to 4.4\% after masked training (\cref{fig:faithfulness}).}
\label{fig:reasoning-examples}
\end{figure*}

